\section{ImageNet-Grid Construction and Its Role in Identifying VIR Heads}
\label{sec:imagenet-grid}

\subsection{Construction}

Every discovery and causal-steering sample in this project is built from
\texttt{vis\_head/imagenet\_grid.py} (\texttt{sample\_grid} and, for the
position-controlled variant, \texttt{sample\_circular\_grid\_series} /
\texttt{build\_imagenet\_circular\_grid\_dataset}). A sample tiles
$R = \text{rows}\times\text{cols}$ single-object ImageNet validation images
into one composite grid image. Concretely, for a grid with $\rho \in
\{0,\dots,R-1\}$ cells:

\begin{enumerate}
  \item \textbf{Class sampling.} $R$ distinct class directories
  $c_0,\dots,c_{R-1}$ are drawn without replacement from the full ImageNet
  validation class list, \texttt{rng.choice(len(class\_dirs), size=n\_cells,
  replace=False)}. Distinctness is enforced at the class level: no two cells
  in the same sample can hold the same category.
  \item \textbf{Image sampling.} For each chosen class $c_\rho$, one image is
  drawn uniformly at random from that class's validation images.
  \item \textbf{Square crop and resize.} Each image is center-cropped to a
  square (side $=\min(w,h)$, so no distortion of object shape) and resized to
  a fixed \texttt{cell\_size} (default $256\times256$), via
  \texttt{\_open\_rgb\_square}.
  \item \textbf{Tiling.} The $R$ cell images are pasted into a row-major
  $\text{rows}\times\text{cols}$ canvas with a fixed pixel gap between cells
  (\texttt{\_assemble\_grid}), and each cell's pixel bounding box is recorded
  exactly (\texttt{\_cell\_bboxes}) — this bounding box is ground truth, not
  inferred.
\end{enumerate}

The result is an \texttt{ImageNetGrid} record carrying, alongside the
composite image itself, the per-cell source images, WordNet IDs, human-readable
class names, and exact pixel bounding boxes — i.e.\ a sample where the
region-to-content mapping is known with certainty by construction, not by any
downstream annotation or detection step.

For the position-controlled \emph{circular} variant used in
Section~\ref{sec:raw-head-selection}'s robustness check, one target class
$c^{*}$ and one fixed set of $R-1$ distractor classes are chosen once, and
then $R$ grids are emitted — one per value of $\rho^{*} \in \{0,\dots,R-1\}$
— with $c^{*}$'s image placed at cell $\rho^{*}$ and the \emph{same}
distractor images filling every other cell in every variant. Content is thus
held fixed while only the target's grid position varies across the $R$-sample
series.

\subsection{Motivation}

A visual information routing (VIR) head is, by hypothesis, an attention head
whose query-to-image attention pattern is causally responsible for pulling a
specific region's visual content into the residual stream at the position
where the model commits to an answer. Locating such heads reliably —
and only such heads — requires an image/prompt distribution with three
properties that natural photographs (e.g.\ COCO) do not offer directly:

\begin{description}
  \item[Known, disjoint ground-truth regions.] The raw-score pipeline
  (Section~\ref{sec:raw-head-selection}) needs a token-to-region map
  $r(t)$ to aggregate attention per region at all. In a natural photo, "where
  is the object" is itself an open detection problem — using a grid instead
  makes region membership a closed-form geometric fact
  (\texttt{assign\_grid\_cells\_to\_tokens} from the recorded \texttt{cell\_bboxes}
  / vision-patch grid, no detector, no annotation noise). Any imprecision in
  which head "found" the object would otherwise be entangled with imprecision
  in what counts as the object's region.
  \item[Referential unambiguity.] Prompts like \texttt{"Identify the
  \{name\}."} are only unambiguous if the named category occurs in exactly
  one cell. Enforcing distinct classes per cell (no-replacement sampling)
  removes the alternative explanation "the head attends to \emph{a} correct
  cell, just not \emph{the} one referenced" — a confound that a natural
  multi-instance photo cannot rule out.
  \item[Independent object/position pairing across samples.] Because
  \texttt{sample\_grid} redraws which class occupies which cell on every
  sample, discovery scores are averaged over many different
  (object, position) pairings. A head that scores well here cannot be a head
  that merely likes one particular object's visual features regardless of
  position, nor a head that merely likes one particular cell regardless of
  content — surviving this average is evidence of genuine
  \emph{content-general, position-conditioned routing}, which is the specific
  behavior "VIR head" is meant to name. The circular-grid variant sharpens
  this further by holding content fixed and varying \emph{only} position
  within a single object's series, directly testing whether a candidate
  head's score is stable under a pure position perturbation (see
  \texttt{score\_variance\_stats} / \texttt{snr\_score}) rather than being an
  artifact of one lucky object/position combination.
\end{description}

Without this construction, a high raw discovery score $\bar{s}_{l,h}$ would
be uninterpretable: it could reflect true region-conditioned routing, or
equally well reflect an object-identity bias, a fixed positional bias, or an
annotation artifact from an imperfect region definition. The grid's controlled
class-distinctness, exact region ground truth, and (in the circular variant)
explicit position/content decoupling are what let a high $\bar{s}_{l,h}$ — and
the resulting causal steering effect once $\mathcal{C}_{\text{raw}}$ is
intervened on — be attributed specifically to visual information routing
rather than to a simpler, non-routing correlate.

\section{Raw-Attention Candidate Head Selection}
\label{sec:raw-head-selection}

This section formalizes the raw-score discovery pipeline exactly as
implemented in \texttt{vis\_head/gaze.py}
(\texttt{collect\_last\_query\_attentions},
\texttt{aggregate\_region\_attention}, \texttt{rank\_heads\_by\_score}) and the
per-checkpoint discovery loops that call them.

\subsection{Attention extraction at the query position}

For a single discovery sample, the multimodal input is a grid image
partitioned into $R = \text{rows}\times\text{cols}$ regions (grid cells) and a
text prompt referencing one target region by object name, e.g.\
\texttt{"Identify the \{name\}."}. One prefill forward pass is run with
\texttt{attn\_implementation="eager"} so per-head attention weights are
materialized. Let $L$ be the number of transformer layers, $H$ the number of
attention heads per layer, and $T$ the prompt length (sequence length at
prefill, before any generation). For layer $l \in \{0,\dots,L-1\}$ and head
$h \in \{0,\dots,H-1\}$, let
\[
  A^{(l,h)} \in \mathbb{R}^{T \times T}
\]
be that head's post-softmax attention matrix from the eager attention
implementation, i.e.\ $\sum_{t=1}^{T} A^{(l,h)}_{q,t} = 1$ for every query
position $q$. Only the attention row at the \emph{last} prompt-token position
$q^{*} = T$ is retained (the position immediately preceding generation,
where the model's committed answer is read off):
\[
  a^{(l,h)}_{t} \;=\; A^{(l,h)}_{q^{*},\, t}, \qquad t = 1,\dots,T .
\]
Stacking over all layers and heads gives the raw query-attention tensor
\[
  \mathbf{a} \in \mathbb{R}^{L \times H \times T},
  \qquad
  \mathbf{a}_{l,h,t} = a^{(l,h)}_{t}.
\]
This is exactly \texttt{collect\_last\_query\_attentions}: one prefill pass,
one attention row (\texttt{attention[0, :, -1, :]}) taken per layer, stacked
into an $(L, H, T)$ array.

\subsection{Region aggregation (no renormalization)}

Let $[t_{\text{img}}^{\text{start}}, t_{\text{img}}^{\text{end}})$ be the
contiguous span of image-token positions in the sequence (architecture-specific,
resolved by \texttt{find\_image\_token\_range}), and let
$r : \{t_{\text{img}}^{\text{start}}, \dots, t_{\text{img}}^{\text{end}}-1\}
\to \{0,\dots,R-1\}$
map each image token to the grid cell it belongs to (\texttt{region\_ids},
produced by \texttt{assign\_grid\_cells\_to\_tokens}). Define the one-hot
region-membership matrix $\mathbf{M} \in \{0,1\}^{T_{\text{img}} \times R}$,
$\mathbf{M}_{t,\rho} = \mathbb{1}[r(t) = \rho]$, where
$T_{\text{img}} = t_{\text{img}}^{\text{end}} - t_{\text{img}}^{\text{start}}$.
The per-region raw score tensor is the attention mass summed over each
region's image tokens:
\[
  \mathbf{s} \in \mathbb{R}^{L \times H \times R},
  \qquad
  \mathbf{s}_{l,h,\rho}
    \;=\; \sum_{t = t_{\text{img}}^{\text{start}}}^{t_{\text{img}}^{\text{end}}-1}
          \mathbf{a}_{l,h,t}\, \mathbf{M}_{t - t_{\text{img}}^{\text{start}},\,\rho} .
\]
Implementation-wise this is the single einsum
\texttt{np.einsum("lht,tr->lhr", image\_attention, region\_onehot)}. Critically,
\emph{no renormalization} is applied across regions: rows of $\mathbf{s}$ do
not sum to $1$. This is a deliberate design choice (documented directly in the
docstring): a head that never attends to image tokens at all, or that spreads
attention diffusely over the whole image, must score near zero rather than
being inflated by a per-row softmax/renormalization step that would make any
head "look diagonal" after the fact. A head only earns a high raw score if it
\emph{both} (a) allocates attention mass to image tokens at all, and (b)
concentrates that mass on the specific queried region.

\subsection{Accumulation over discovery samples}

For a discovery run of $N$ i.i.d.\ grid samples, sample $n$ has a
ground-truth target region $\rho^{*}_n$ (the grid cell holding the object
named in that sample's prompt). Only the score column at the \emph{true}
target region is kept per sample:
\[
  \sigma^{(n)}_{l,h} \;=\; \mathbf{s}^{(n)}_{l,h,\rho^{*}_n}.
\]
The raw discovery score for head $(l,h)$ is the sample mean over all
successfully processed samples ($\mathcal{N} \subseteq \{1,\dots,N\}$, valid
samples after discarding any that raised an exception during forward/region
assignment):
\[
  \boxed{\;
  \bar{s}_{l,h} \;=\; \frac{1}{|\mathcal{N}|}\sum_{n \in \mathcal{N}} \sigma^{(n)}_{l,h}
  \;}
\]
i.e.\ \texttt{vis\_head\_scores = raw\_sum / valid} where
\texttt{raw\_sum += region\_attn[:, :, target\_cell]} is accumulated
incrementally inside the discovery loop. $\bar{\mathbf{s}} \in \mathbb{R}^{L
\times H}$ is the final raw score matrix passed to selection.

\subsection{Top-$K$ candidate selection}

Candidate heads are chosen by a pure descending sort of $\bar{\mathbf{s}}$,
with \emph{no normalization of any kind} applied before ranking
(\texttt{rank\_heads\_by\_score}):
\[
  \text{rank}(l,h) \;=\; \bigl|\{(l',h') : \bar{s}_{l',h'} > \bar{s}_{l,h}\}\bigr| + 1,
\]
\[
  \mathcal{C}_{\text{raw}}
    \;=\; \operatorname*{arg\,top\text{-}K}_{(l,h) \,\in\, \{0,\dots,L-1\}\times\{0,\dots,H-1\}} \bar{s}_{l,h}
    \;=\; \{(l,h) : \text{rank}(l,h) \le K\}.
\]
$\mathcal{C}_{\text{raw}}$ (size $K$, $K{=}15$ throughout this project) is
the raw-attention candidate head set. It is exactly this set (or a
normalized-score variant of it, e.g.\ per-layer max/z-score-normalized
$\bar{\mathbf{s}}$ ranked the same way) that is grouped by layer
(\texttt{group\_heads\_by\_layer}) and handed to the causal steering
intervention (\texttt{make\_static\_attention\_mask\_hook} /
\texttt{register\_mask\_hooks}), which forcibly boosts attention to the
target region and suppresses attention elsewhere, restricted to exactly the
heads in $\mathcal{C}_{\text{raw}}$, to test whether the \emph{observationally}
selected heads are also \emph{causally} responsible for the model's behavior.

\subsection{Summary}

\[
\underbrace{A^{(l,h)}}_{\text{eager attn.\ matrix}}
\;\xrightarrow{\text{last query row}}\;
\underbrace{a^{(l,h)}_t}_{\mathbf{a}\in\mathbb{R}^{L\times H\times T}}
\;\xrightarrow{\text{region sum, unnormalized}}\;
\underbrace{\mathbf{s}_{l,h,\rho}}_{\text{per-sample, per-region}}
\;\xrightarrow{\text{select }\rho=\rho^{*},\ \text{mean over }N\text{ samples}}\;
\underbrace{\bar{s}_{l,h}}_{\text{raw discovery score}}
\;\xrightarrow{\text{descending sort, top-}K}\;
\mathcal{C}_{\text{raw}}.
\]
No step in this pipeline divides by a per-layer or per-head normalizer;
normalization variants (per-layer max, per-layer z-score, SNR across samples)
are computed only as \emph{alternative} rankings from the same
$\bar{\mathbf{s}}$ (or per-sample $\sigma^{(n)}_{l,h}$) tensor, for direct
comparison against $\mathcal{C}_{\text{raw}}$.

\section{Causal Evaluation: Intervention, MCQ Generation, and Probing}
\label{sec:causal-eval}

A candidate head set $\mathcal{C}$ (e.g.\ $\mathcal{C}_{\text{raw}}$ from
Section~\ref{sec:raw-head-selection}, or any normalized-score variant) is a
purely \emph{observational} object until it is shown to be
\emph{causally} load-bearing: forcibly redirecting exactly those heads'
attention must change the model's behavior in the predicted direction. This
section formalizes the two independent measurement instruments used
throughout the project to test that -- one behavioral (grid MCQ generation
with ground-truth-scored predictions), one representational (a linear probe
on hidden states) -- both applied to the same held-out sample set under the
same intervention, so their outcomes can be directly cross-checked against
one another.

\subsection{Causal intervention: static attention-mask biasing}

Given target region $\rho^{*}$ and its image-token positions
$T_{\rho^{*}}$, and the remaining image-token positions
$T_{\text{other}} = \{t_{\text{img}}^{\text{start}},\dots,t_{\text{img}}^{\text{end}}-1\} \setminus T_{\rho^{*}}$,
the \texttt{"boost\_suppress"} intervention mode
(\texttt{intervention\_positions}) designates
\[
  \text{boost positions} = T_{\rho^{*}}, \qquad
  \text{suppress positions} = T_{\text{other}},
\]
leaving all non-image (text/prompt) positions untouched. For every candidate
head $(l,h) \in \mathcal{C}$, a static additive bias is injected into that
head's pre-softmax attention logits at \emph{every} decoding step
(\texttt{make\_static\_attention\_mask\_hook}, attached via a
\texttt{register\_forward\_pre\_hook} on \texttt{self\_attn}, i.e.\ before
the architecture's own attention-mask combination):
\[
  \widetilde{\text{mask}}^{(l,h)}_{1,t}
  \;=\;
  \text{mask}_{1,t}
  \;+\;
  \begin{cases}
    +\beta, & t \in T_{\rho^{*}} \text{ and } h \in \mathcal{C}_l \\
    -\beta, & t \in T_{\text{other}} \text{ and } h \in \mathcal{C}_l \\
    0, & \text{otherwise},
  \end{cases}
\]
where $\mathcal{C}_l = \{h : (l,h) \in \mathcal{C}\}$ is the set of
intervened heads in layer $l$, and $\beta = \texttt{DEFAULT\_SWAP\_BIAS} =
10{,}000$ -- large enough that, after softmax, attention mass on suppressed
positions is driven to numerical zero and mass on boosted positions
dominates, regardless of the head's original (pre-intervention) attention
pattern. Heads not in $\mathcal{C}$, and all layers not containing an
intervened head, are left completely unmodified. Only heads in
$\mathcal{C}$ at layer $l$ receive the bias; every other head at that same
layer keeps its original attention. This mask is added identically at every
generation step (prefill and every decode step alike, \texttt{decode\_only=
False} in this project's usage), so the boosted/suppressed routing is
enforced for the entire response, not just the first token.

\subsection{Instrument 1: MCQ generation with ground-truth-scored prediction}

For a $k$-way multiple-choice sample with option letters
$\mathcal{L} = \{A,\dots\}$ ($|\mathcal{L}|=k$, $k{=}4$ throughout) and
correct letter $\ell^{*} \in \mathcal{L}$, the model is run under
\texttt{model.generate(..., do\_sample=False, output\_scores=True,
return\_dict\_in\_generate=True)} (greedy decoding) with the intervention
hooks of the previous subsection registered for the duration of generation.
Let $g_1,\dots,g_M$ be the generated token ids (excluding the prompt) and
$\text{ids}(\ell)$ the (possibly multi-token, whitespace-variant) token-id
set naming letter $\ell$ (\texttt{get\_letter\_token\_ids}). The
\emph{answer step} is the position of the \emph{last} generated token that
names any option letter,
\[
  m^{*} \;=\; \max\bigl\{\, m \in \{1,\dots,M\} \;:\; g_m \in \textstyle\bigcup_{\ell \in \mathcal{L}} \text{ids}(\ell) \,\bigr\},
\]
taking the last rather than the first such occurrence so that both
instant-answer models (letter is $g_1$) and chain-of-thought models (the
letter may recur inside intermediate reasoning tokens before the final
commitment) are scored at the model's actual final decision point. The
predicted letter is read off directly from the generated token,
$\hat\ell = \ell$ such that $g_{m^{*}} \in \text{ids}(\ell)$, and a genuine
probability vector is recovered from the real logits at that exact decode
step (not from a softmax over the whole vocabulary, and not requiring a
second forward pass):
\[
  p_\ell \;=\; \operatorname{softmax}_{\ell \in \mathcal{L}}
  \Bigl(\max_{v \,\in\, \text{ids}(\ell)} z^{(m^{*})}_v\Bigr),
  \qquad
  z^{(m^{*})} = \texttt{out.scores}[m^{*}][0],
\]
i.e.\ the max logit over that letter's token-id variants at the answer
step, softmax-normalized over the $k$ options. Aggregating $\hat\ell$ and
$p$ over $N$ held-out samples with ground truth $\ell^{*}_n$ gives standard
classification metrics -- accuracy $\frac{1}{N}\sum_n
\mathbb{1}[\hat\ell_n = \ell^{*}_n]$, macro-F1, and macro one-vs-rest AUC
from $p$ -- computed once with no intervention (\emph{baseline}, $\mathcal{C}
= \varnothing$) and once with the candidate heads intervened
(\emph{steered}). The behavioral causal effect of $\mathcal{C}$ is the
paired difference $\text{acc}_{\text{steered}} - \text{acc}_{\text{baseline}}$
on the identical $N$ samples, with significance assessed by McNemar's test
on the paired correct/incorrect outcome vectors (2x2 discordant-pair count,
$\chi^2_1$ with continuity correction).

\subsection{Instrument 2: linear-probe classification on hidden states}

The generation-based measurement above is behavioral: it can only observe
whether the sampled/greedy output token sequence changed. A second,
purely \emph{representational} measurement asks whether the model's internal
state at the decision point was linearly separable by the correct answer,
independent of whatever the decoding procedure ultimately emits. For the
identical sample and identical intervention (hooks registered or not), one
extra forward pass with \texttt{output\_hidden\_states=True} is run (no
generation), and the last transformer layer's hidden state at the final
prompt token -- the representation immediately before the model would
start generating -- is extracted:
\[
  \mathbf{h}_n \;=\; H^{(L)}_n[0,\,-1,\,:] \;\in\; \mathbb{R}^{d_{\text{model}}},
  \qquad H^{(L)}_n = \texttt{out.hidden\_states[-1]}.
\]
(For architectures whose final hidden-state tensor carries an extra
leading dimension beyond the standard $(\text{batch}, \text{seq},
\text{hidden})$ -- e.g.\ Gemma-3n's AltUp, which stacks
\texttt{config.altup\_num\_inputs} parallel prediction streams as
$(\text{streams}, \text{batch}, \text{seq}, \text{hidden})$ -- the stream at
\texttt{config.altup\_active\_idx} is selected first, since that is the
stream architecturally used to produce the output logits; the standard
3-D case is unaffected.)

Stacking $\mathbf{h}_1,\dots,\mathbf{h}_N$ with their ground-truth labels
$\ell^{*}_1,\dots,\ell^{*}_N$ gives a $k$-class classification dataset. A
multinomial logistic-regression probe is fit on a $70/30$ train/test split
(\texttt{fit\_and\_eval\_probe}, \texttt{sklearn.linear\_model.
LogisticRegression}, $C{=}1$):
\[
  \hat{\ell}\,(\mathbf{h}) \;=\; \operatorname*{arg\,max}_{\ell \in \mathcal{L}}
  \; \mathbf{w}_\ell^{\top}\mathbf{h} + b_\ell,
  \qquad
  (\mathbf{w}_{\bullet}, \mathbf{b}_{\bullet}) = \operatorname*{arg\,min}
  \sum_{n \in \text{train}} -\log \operatorname{softmax}(\mathbf{W}\mathbf{h}_n+\mathbf{b})_{\ell^{*}_n}
  \;+\; \tfrac{1}{2C}\lVert \mathbf{W}\rVert_F^2 ,
\]
reporting held-out accuracy and macro one-vs-rest AUC. This probe is
re-fit independently for the baseline hidden states $\{\mathbf{h}^{\text{base}}_n\}$
and the steered hidden states $\{\mathbf{h}^{\text{steer}}_n\}$ (same $N$
samples, same train/test split indices), and the \emph{probe ATE} is the
accuracy (or AUC) difference between the two fits. Because this measures
linear separability of a representation rather than a decoded token, it can
detect a causal shift that generation-based scoring misses (the correct
answer becomes decodable from $\mathbf{h}$ but the model's greedy decoding
still emits something else) or, symmetrically, confirm that a generation-level
effect is backed by an actual representational change rather than a
decoding-surface artifact.

\subsection{Combined interpretation}

The two instruments are run on the \emph{same} $N$ samples under the
\emph{same} baseline/steered conditions, giving four paired accuracy
numbers per checkpoint per candidate set $\mathcal{C}$:
$(\text{acc}^{\text{gen}}_{\text{base}}, \text{acc}^{\text{gen}}_{\text{steer}},
\text{acc}^{\text{probe}}_{\text{base}}, \text{acc}^{\text{probe}}_{\text{steer}})$.
Agreement between the generation ATE and the probe ATE is the strongest
form of evidence that $\mathcal{C}$ is causally load-bearing for the task
in the way intended; a probe-only effect with no generation-level effect
localizes the causal claim to representation rather than behavior, and a
generation-only effect with no probe shift is grounds to suspect an
artifact of the decoding procedure (e.g.\ the forced attention mask
perturbing which token is sampled without genuinely changing what
information is linearly available) rather than of the underlying
computation.
